\section{The one-patch calculation}
\label{app:patch-calculations}

\begin{discussion}
  \item State that this appendix contains the local calculation used in
  the representation, positivity criterion, critical-density formula,
  pure-death comparison, and end-factor relaxation.
\end{discussion}

\begin{notation}
\label{not:patch-calculation}
Fix a patch \(P\), write
\[
  i=i(P),\qquad s=s(P),\qquad e=e(P),\qquad S=S(P),
\]
and let \(s\le t\le e\). Put
\[
  \Delta_-=t-s,
  \qquad
  \Delta_+=e-t,
  \qquad
  \Delta=e-s.
\]
Define
\begin{align}
  \alpha_i
  &=
  \sum_{R\subseteq N(i)}\delta_i(R)
  +
  \sum_{\vn\ne R\subseteq N(i)}\beta_i(R),
  \label{eq:alpha-i}
  \\
  V_i
  &=
  \alpha_i+a_i^\beta(\vn),
  \label{eq:Vi-alpha}
  \\
  \varphi_i(\Delta)
  &=
  e^{-\alpha_i\Delta}
  +
  \delta_i(\vn)
  \int_0^\Delta e^{-\alpha_iw}\,\dd w,
  \label{eq:varphi}
  \\
  \psi_i(\Delta_-,\Delta_+,z)
  &=
  \delta_i(\vn)
  \int_0^{\Delta_-}e^{a_i^\beta(\vn)w}\,\dd w
  +
  z e^{a_i^\beta(\vn)\Delta_-}
  \varphi_i(\Delta_+).
  \label{eq:psi-two-time}
\end{align}
Write \(\psi_i(\Delta,z)=\psi_i(\Delta,0,z)\).
\end{notation}

\begin{proposition}[Explicit patch factors]
\label{prop:explicit-patch-factors}
For \(s\le t\le e<\infty\), the patch contribution up to time \(t\)
is
\begin{equation}
  C_t(z,P)
  =
  \begin{cases}
  \dfrac{
    \psi_i(\Delta_-,\Delta_+,z)
  }{
    \varphi_i(\Delta)
  },
  &
  \mathsf X(P)=\sI,\quad
  \mathsf Y(P)\in\{\sI,\sE\},
  \\[1.2em]
  z e^{V_i\Delta_-},
  &
  (\mathsf X(P),\mathsf Y(P))=(\sI,\sO),
  \\[1.2em]
  \dfrac{
    \delta_i(S)\sigma_i^\delta(S)
    +
    \beta_i(S)\sigma_i^\beta(S)
    \psi_i(\Delta_-,\Delta_+,z)
  }{
    \delta_i(S)+\beta_i(S)\varphi_i(\Delta)
  },
  &
  \mathsf X(P)=\sO,\quad
  \mathsf Y(P)\in\{\sI,\sE\},
  \\[1.5em]
  \sigma_i^\beta(S)z e^{V_i\Delta_-},
  &
  (\mathsf X(P),\mathsf Y(P))=(\sO,\sO).
  \end{cases}
  \label{eq:explicit-patch-factors}
\end{equation}
For a full patch,
\[
  C(P)=\lim_{t\uparrow e(P)}C_t(1,P).
\]
The same endpoint definition applies when \(e(P)=\infty\).
\end{proposition}

\begin{proof}
If the patch starts active and has incoming or end terminal type,
consistency occurs if no outgoing interaction happens before \(e\), or
if the first one is an empty-target death. Thus
\begin{equation}
  \pP_P(\operatorname{Con}(P))
  =
  e^{-\alpha_i\Delta}
  +
  \delta_i(\vn)
  \int_0^\Delta e^{-\alpha_iw}\,\dd w
  =
  \varphi_i(\Delta).
  \label{eq:I-normalizer}
\end{equation}
Splitting according to whether that death occurs before or after
\(t\), the unnormalized weighted expectation is
\begin{align*}
  &\delta_i(\vn)
  \int_0^{\Delta_-}e^{(V_i-\alpha_i)w}\,\dd w
  \\
  &\quad+
  z e^{(V_i-\alpha_i)\Delta_-}
  \left[
    e^{-\alpha_i\Delta_+}
    +
    \delta_i(\vn)
    \int_0^{\Delta_+}e^{-\alpha_iw}\,\dd w
  \right]
  =
  \psi_i(\Delta_-,\Delta_+,z).
\end{align*}
Division by \eqref{eq:I-normalizer} gives the first row.

If the patch begins at an outgoing interaction, its initial mark is a
split or birth with probabilities proportional to
\(\delta_i(S)\) and \(\beta_i(S)\). A split makes the source inactive,
whereas a birth leaves it active. For incoming or end terminal type,
the consistency probability and unnormalized signed weight are,
respectively,
\[
  \frac{
    \delta_i(S)+\beta_i(S)\varphi_i(\Delta)
  }{
    \delta_i(S)+\beta_i(S)
  }
\]
and
\[
  \frac{
    \delta_i(S)\sigma_i^\delta(S)
    +
    \beta_i(S)\sigma_i^\beta(S)
    \psi_i(\Delta_-,\Delta_+,z)
  }{
    \delta_i(S)+\beta_i(S)
  }.
\]
Their ratio is the third row. Outgoing terminal consistency forces the
source to remain active throughout the patch; for an outgoing-initial
patch it also forces the initial mark to be a birth. This gives the
second and fourth rows.
\end{proof}

\begin{corollary}[Spin-rate numerator]
\label{cor:spin-rate-numerator}
The numerator in the outgoing-initial rows is
\begin{equation}
  c_i^0(S)
  -
  \bigl(c_i^0(S)+c_i^1(S)\bigr)
  \psi_i(\Delta_-,\Delta_+,z).
  \label{eq:OI-spin-numerator}
\end{equation}
The outgoing-terminal sign is nonnegative precisely when
\[
  c_i^0(S)+c_i^1(S)\le0.
\]
\end{corollary}

\begin{proof}
Use
\[
  \delta_i(S)\sigma_i^\delta(S)=c_i^0(S),
  \qquad
  \beta_i(S)\sigma_i^\beta(S)
  =-c_i^0(S)-c_i^1(S)
\]
in \eqref{eq:explicit-patch-factors}.
\end{proof}

\begin{proposition}[No-interaction continuation]
\label{prop:continuation-relaxation}
Put \(r_i=c_i^0(\vn)+c_i^1(\vn)\). If \(r_i>0\), define
\[
  q_i=\frac{c_i^0(\vn)}{r_i}.
\]
Then
\begin{equation}
  \psi_i(u,z)
  =
  q_i+(z-q_i)e^{-r_iu},
  \qquad
  \psi_i(u+v,z)
  =
  \psi_i\bigl(u,\psi_i(v,z)\bigr).
  \label{eq:empty-neighbour-relaxation}
\end{equation}
If \(r_i=0\), then \(\psi_i(u,z)=z\).

Let \(Q\) be an end patch based at \(i\) at time \(T\). Define its
continuation through \(u\ge T\) without a successful interaction by
\begin{equation}
  C(z,Q^{\uparrow u})
  :=
  C\bigl(\psi_i(u-T,z),Q\bigr).
  \label{eq:continuation-contribution}
\end{equation}
If \(C(z,Q)=a(Q)+b(Q)z\), then
\begin{equation}
  \partial_zC(z,Q^{\uparrow u})
  =
  e^{-r_i(u-T)}b(Q).
  \label{eq:continuation-slope}
\end{equation}
For \(r_i>0\),
\[
  C(Q^{\uparrow\infty})
  =
  \lim_{u\to\infty}C(z,Q^{\uparrow u})
  =
  C(q_i,Q).
\]
\end{proposition}

\begin{proof}
Substitute \(a_i^\beta(\vn)=-r_i\) and
\(\delta_i(\vn)=c_i^0(\vn)\) into
\eqref{eq:psi-two-time}. The remaining assertions follow from the
affinity of \(C(z,Q)\).
\end{proof}
